
\chapter{Análise do Código MATLAB: \texttt{binrelsanalyser.m}}
\author{Expert Software Engineer}
\date{\today}
\section{Visão Geral Inicial}

\textbf{What is the main purpose of this file?} \\
The main purpose of this MATLAB function is to analyze a binary relation (represented by an adjacency matrix) and classify it based on its fundamental mathematical properties: reflexivity, transitivity, symmetry, and anti-symmetry. Additionally, it computes corrective transformations, generating a transitive closure if the relation lacks transitivity, or condensing it into a partial order (quotient set) if it is a preorder lacking anti-symmetry.

\textbf{What type of analysis or calculation does it perform?} \\
It performs matrix-based logical analysis on a sparse adjacency matrix. The calculations rely heavily on matrix multiplication (to test path connectivity), logical indexing, matrix transposition, and identity matrix comparisons to test structural relation rules.

\textbf{What is the mathematical/theoretical context?} \\
This file operates within the realms of \textbf{Grafo Theory} and \textbf{Order Theory}. It mathematically evaluates binary relations ($R \subseteq V \times V$), identifying specific structures like \textit{preorders} (reflexive and transitive), \textit{equivalence relations} (preorders that are symmetric), and \textit{partial orders} (preorders that are anti-symmetric). The transformation to \texttt{Ranti} implicitly uses the category theory concept of a \textit{coequalizer} to build a quotient space $R/{\sim}$.

\vspace{1em}

\vspace{1em}

\section{Análise Passo a Passo do Código}

\subsection{1. Título da Secção/Função: Variable Initialization and Squareness Check}
\textbf{Explanation:} \\
This section initializes the default property output (\texttt{prop = 6}, representing an unclassified state) and empty arrays for the transformed matrices. It then checks if the input matrix $R$ is square ($n_1 = n_2$). A binary relation on a single set $V$ must map $V \to V$, requiring a square adjacency matrix. If it is not square, the function alerts the user, sets \texttt{prop = 0}, and terminates early.

\begin{lstlisting}[language=Matlab]
prop=6; Rtran=[]; Ranti=[];
[n1,n2]=size(R); idV=(1:n1)'; 

if ~isequal(n1,n2),
    disp('R is not a square matrix')
    prop=0;
    return
end
\end{lstlisting}

\subsection{2. Título da Secção/Função: Reflexivity Check}
\textbf{Explanation:} \\
This block evaluates whether the relation is reflexive. A relation is reflexive if every element relates to itself ($\forall x \in V, xRx$). In an adjacency matrix, this means the main diagonal must consist entirely of ones. The code compares the extracted diagonal \texttt{diag(R)} with a column vector of ones. It assigns \texttt{prop = 1} if true.

\begin{lstlisting}[language=Matlab]
if ~isequal(diag(R),ones(n1,1))
    disp('R is not relexive')
    prop=0;
else
    disp('R is reflexive')
    prop=1;
end
\end{lstlisting}

\subsection{3. Título da Secção/Função: Transitivity Check and Transitive Closure Computation}
\textbf{Explanation:} \\
This section checks for transitivity. Mathematically, transitivity dictates that if $aRb$ and $bRc$, then $aRc$. In matrix terms, the square of the adjacency matrix ($R^2$) must be a subset of $R$ (\texttt{logical(R\textasciicircum 2) <= R}). If the relation is reflexive but not transitive, it computes the \textit{transitive closure} ($R^+$) using a loop that iteratively adds connections (\texttt{Rn + Rn*R}) until the matrix stops changing (a variant of the Warshall algorithm). If it is already transitive and reflexive, it is classified as a \textit{preorder} (\texttt{prop = 2}).

\vspace{1em}

\vspace{1em}

\begin{lstlisting}[language=Matlab]
if ~(all(all(logical(R^2)<=R))) && prop==1
    disp('R is not transitive, the output Rtran is its transitive closure')
    Rn=R; 
    Rtran=logical(Rn+Rn*R);
    while ~isequal(Rn,Rtran)
        Rn=Rtran;
        Rtran=logical(Rn+Rn*R);
    end    
elseif prop==1
    prop=2;
    disp('R is reflexive and transitive (also called a preorder)')
end
\end{lstlisting}

\subsection{4. Título da Secção/Função: Symmetry Check and Equivalence Classification}
\textbf{Explanation:} \\
Here, the code checks if the relation is symmetric by comparing the matrix $R$ to its transpose $R'$ (\texttt{isequal(R, R')}). If it is symmetric, and previously established as a preorder (\texttt{prop == 2}), it is classified as an \textit{equivalence relation} (\texttt{prop = 3}). 

\begin{lstlisting}[language=Matlab]
if isequal(R,R')
    disp('R is symmetric')
    if prop==2, prop=3; disp('R is an equivalence relation') end
\end{lstlisting}

\subsection{5. Título da Secção/Função: Anti-Symmetry Check and Partial Order Condensation}
\textbf{Explanation:} \\
If the relation is not symmetric, it tests for anti-symmetry. An anti-symmetric matrix has no overlapping 1s across its main diagonal, except on the diagonal itself (\texttt{R \& R' == eye}). 
If the relation is neither symmetric nor anti-symmetric, but is a preorder, the code "forces" anti-symmetry. It finds bidirectional connections (\texttt{find(R \& R')}), uses a \texttt{coequalizer} function to define equivalent classes ($x \sim y$), and constructs a quotient space matrix (\texttt{Ranti}), effectively converting the preorder into a \textit{partial order}. If it was natively anti-symmetric and a preorder, it is classified as a partial order (\texttt{prop = 4}).

\begin{lstlisting}[language=Matlab]
elseif ~isequal((R & R'),eye(size(R)))
    disp('R is not symmetric (hence not an equivalence relation) nor anti-symmetric (hence not a partial order)')
    if prop==2 
        prop=5;
        disp('but it is a preorder (reflexive and transitive), so, forcing anti-symmetry by identifying x~y iff R(x,y) and R(y,x), a partial order is obtained,')
        disp('Ranti is the adjacency matrix of the partial order R/~')
        [s,t]=find(R & R');   
        q=coequalizer(s,t);  
        Ranti=R(q==idV,q==idV);
    end
else
    disp('R is anti-symmetric')
    if prop==2, prop=4; disp('R is a partial order'), end
end
\end{lstlisting}

\section{Saídas e Elementos Visuais Esperados}

Although this specific script does not natively call \texttt{plot} or \texttt{digraph} functions to draw figures on the screen, its mathematical outputs map directly to standard visual representations in graph theory.

\textbf{Console Saídas:} \\
The script primarily prints strings to the MATLAB Command Window tracking the logical flow. Expect statements like \texttt{'R is reflexive and transitive (also called a preorder)'} or \texttt{'R is not transitive, the output Rtran is its transitive closure'}, serving as a real-time diagnostic log.

\textbf{Variable Saídas translated to Grafos:}
\begin{itemize}
    \item \textbf{Initial Adjacency Matriz ($R$):} Translates to a standard directed graph (digraph). The vertices (axes/nodes) represent the elements $1$ to $n_1$. An edge exists from node $i$ to $j$ if $R(i,j) = 1$.
    \item \textbf{Transitive Closure ($R_{tran}$):} Visually, this graph looks identical to the original graph $R$, but with all possible "shortcut" edges drawn. If there is a path from node $A$ to node $C$ through node $B$, $R_{tran}$ visually adds a direct edge from $A \to C$.
    \item \textbf{Quotient Partial Order ($R_{anti}$):} If you were to plot this, you would see a \textit{Directed Acyclic Grafo (DAG)} or a \textit{Hasse Diagram}. Visually, clusters of nodes that previously had 2-way arrows (cycles indicating $A \to B$ and $B \to A$) are condensed/collapsed into a single representative "super-node" by the coequalizer. The axes of this new matrix no longer represent the original vertices, but rather the equivalence classes (partitions) of the vertices.
\end{itemize}

\vspace{1em}


